\section*{Aktivitas Pembelajaran dan Praktikum}
\addcontentsline{toc}{section}{Aktivitas Pembelajaran dan Praktikum}

\begin{table}[H]
\centering
\caption{Rencana aktivitas Bab XIV}
\label{tab:ch14-aktivitas}
\begin{tabularx}{\textwidth}{@{}llX@{}}
\toprule
\textbf{Sesi} & \textbf{Bentuk} & \textbf{Kegiatan} \\
\midrule
1 & Ceramah & Ruang lingkup analisis semantik \\
2 & Latihan & Menulis tabel aturan semantik ekspresi dan pernyataan \\
3 & Diskusi & Pengujian aturan pada pohon sintaks \\
4 & Praktikum & P4-A: tabel aturan, uji pohon, dan inferensi tipe \\
\bottomrule
\end{tabularx}
\end{table}

\subsection*{P4-A: Tabel Aturan Semantik dan Inferensi Tipe}
\textbf{Tujuan.} Membuat tabel aturan semantik untuk ekspresi dan
pernyataan, mengujinya pada pohon sintaks, lalu
mengimplementasikan inferensi tipe.

\textbf{Langkah kerja.}
\begin{enumerate}[label=\arabic*.,leftmargin=*]
  \item Lengkapi tabel aturan semantik ekspresi pada
        \tabref{tab:ch14-sdd} dan tabel pernyataan pada
        \tabref{tab:ch14-sdd-stmt}.
  \item Pilih satu program yang memuat deklarasi, assignment, dan
        percabangan atau perulangan; gambar pohon sintaksnya.
  \item Terapkan setiap aturan pada simpul pohon dan catat hasil atau
        \code{TYPE\_ERROR}.
  \item Lengkapi \file{include/types.h} dan \file{src/types.c}.
  \item Implementasikan \code{infer\_expr} dan \code{check\_stmt}.
  \item Uji minimal 15 ekspresi dan 5 pernyataan, termasuk campuran
        tipe dan identifier yang tidak dideklarasikan.
  \item Kumpulkan tabel, pohon teranotasi, dan keluaran pengujian.
\end{enumerate}

\textbf{Definisi selesai.}
\begin{checklist}
  \item Tabel aturan ekspresi dan pernyataan lengkap.
  \item Pohon sintaks teranotasi dengan tipe atau error.
  \item Minimal 15 ekspresi dan 5 pernyataan diuji.
  \item \code{TYPE\_ERROR} dipropagasikan tanpa pesan berulang.
  \item Laporan dan keluaran pengujian diserahkan.
\end{checklist}

P4-A menjadi bagian Praktikum 4 (12\%) yang dilaporkan pada Bab XV.
